\subsection{Winner-take-all collapse under GRPO, Dr.GRPO, and MaxRL}
\label{app:theory-collapse}

Prior work studies objective-induced outcome collapse and inverse-probability
correction \citep{sinha2026expected}; UCPO analyzes correctness-indifferent collapse
through sampling bias and gradient amplification \citep[Section~4]{lochab2026ucpo}.
Here we connect our group estimators to a standard categorical flow and then study
persistent replay over execution-defined keys.
Consider the infinitesimal expected-update flow
$\dot z=\E[\widehat g_G]=c_G(P)\nabla_zP$. This is the clean mean-field limit
of fresh on-policy groups and the first-order PPO update. By Lemmas~\ref{lem:grpo-mean} and~\ref{lem:maxrl-mean}, it covers GRPO, Dr.GRPO, and binary MaxRL with their respective positive functions $c_G$. It deliberately
removes sampling noise; the result below therefore identifies a structural
instability rather than blaming finite batches.

The conditional correct-mode flow derived below is a standard
frequency-dependent replicator system with fitness $f_c(q)=q_c$
\citep{hofbauer1998evolutionary}; the following
theorem is not claimed as a new replicator-dynamics result. Its role is to show
that the expected GRPO, Dr.GRPO, and binary-MaxRL updates reduce to that
familiar winner-take-all flow under the paper's categorical assumptions.

\paragraph{A local view of the instability.}
For $q=(0.60,0.30,0.10)$ in \hyperref[ex:theory-running]{the running example},
$\sum_cq_c^2=0.46$. In the effective time $\tau$ defined in the proof,
Equation~\eqref{eq:grpo-replicator} gives
\[
 \frac{dq}{d\tau}=(0.084,-0.048,-0.036),\qquad
 \frac{d}{d\tau}\log\frac{q_1}{q_2}=0.30.
\]
The two less common modes lose their share of successful outputs even while
correctness increases. This calculation is local; the theorem establishes
that the imbalance persists all the way to the single-mode limit.

\begin{theorem}[Standard replicator specialization for binary-reward RL]
\label{thm:grpo-collapse}
Start from finite logits with $0<P(0)<1$. If the conditional correct-mode
distribution $q(0)$ has a unique largest coordinate $c^\star$, then
\[
  P(t)\longrightarrow1,
  \qquad
  q_{c^\star}(t)\longrightarrow1,
  \qquad
  q_c(t)\longrightarrow0\quad(c\ne c^\star).
\]
Thus these GRPO, Dr.GRPO, and binary MaxRL mean flows converge to a single
correct execution mode for generic initial conditions, even though every mode
in $\mathcal C$ has exactly the same reward.
\end{theorem}

\begin{proof}
For a correct mode $c$ and incorrect category $a$,
\[
  \dot z_c=c_G(P)p_c(1-P),
  \qquad
  \dot z_a=-c_G(P)p_aP.
\]
Also $\dot P=c_G(P)\lVert\nabla_zP\rVert_2^2>0$. If $P$ converged to an
interior value, continuity and positivity of $c_G$, together with
\[
 \lVert\nabla_zP\rVert_2^2
 \ge P^2(1-P)^2\left(\frac1m+\frac1n\right),
\]
would keep $\dot P$ bounded away from zero, a contradiction. Hence $P\to1$.

Define the effective optimization time
\[
  \tau(t)=\int_0^t c_G(P(s))P(s)(1-P(s))\,ds.
\]
If $\tau(\infty)<\infty$, every correct and incorrect logit would have finite
total variation, because the absolute velocity of each is at most $\dot\tau$.
All logits would then converge to finite values, contradicting $P\to1$.
Therefore $\tau(t)\to\infty$, leaving infinite effective time for the
correct-simplex dynamics below.

On the correct simplex, changing from $t$ to $\tau$ gives the replicator flow
\begin{equation}
  \frac{dq_c}{d\tau}
  =q_c\left(q_c-\sum_{d\in\mathcal C}q_d^2\right),
  \qquad
  \frac{d}{d\tau}\log\frac{q_c}{q_d}=q_c-q_d.
  \label{eq:grpo-replicator}
\end{equation}
The initially unique maximum remains the unique maximum. For $c\ne c^\star$,
set $\xi_c=q_c/q_{c^\star}$, so $\xi_c(0)<1$ and $\xi_c$ is nonincreasing.
Since $q_{c^\star}\ge1/m$,
\[
 \frac{d\log\xi_c}{d\tau}
 =-q_{c^\star}(1-\xi_c)
 \le-\frac{1-\xi_c(0)}{m}.
\]
Hence $\xi_c(\tau)\le\xi_c(0)\exp[-(1-\xi_c(0))\tau/m]\to0$.
Every competing mode vanishes relative to $c^\star$; normalization then
yields $q_{c^\star}\to1$, establishing the claimed single-mode limit.
\end{proof}

\begin{corollary}[Ties are nongeneric]
\label{cor:grpo-ties}
If exactly $k<m$ correct modes tie for the largest initial probability, symmetry
preserves the tie and the flow converges to the uniform distribution on those
$k$ modes while eliminating the rest. Only the exactly uniform initialization preserves all $m$ correct modes
in the limit; a tie among $1<k<m$ maxima yields partial collapse onto those
$k$ modes. Any perturbation creating a unique maximum gives single-mode
collapse by Theorem~\ref{thm:grpo-collapse}.
\end{corollary}

\paragraph{An exact tie is a different initial condition.}
Starting instead from $q=(0.45,0.45,0.10)$ gives the limit
$(0.50,0.50,0)$. Changing the initialization to $(0.451,0.449,0.10)$ makes
the first mode the unique winner. Thus an exactly symmetric two-mode limit
does not imply robustness to a small imbalance in the initial policy.

Finite groups also introduce sampling effects: a mode absent from
a group has no sampled response of its own carrying a score term,
all-correct groups have zero task advantage, and sampling noise can break exact
ties. Softmax normalization and shared parameters can still change an
unsampled mode's probability. These facts are not needed for
the theorem.

